
\documentclass{article}
\usepackage{colortbl}
\usepackage{makecell}
\usepackage{multirow}
\usepackage{supertabular}

\begin{document}

\newcounter{utterance}

\centering \large Interaction Transcript for game `codenames', experiment `risk\_low', episode 2 with deepseek{-}v4{-}pro.
\vspace{24pt}

{ \footnotesize  \setcounter{utterance}{1}
\setlength{\tabcolsep}{0pt}
\begin{supertabular}{c@{$\;$}|p{.15\linewidth}@{}p{.15\linewidth}p{.15\linewidth}p{.15\linewidth}p{.15\linewidth}p{.15\linewidth}}
    \# & $\;$A & \multicolumn{4}{c}{Game Master} & $\;\:$B\\
    \hline

    \theutterance \stepcounter{utterance}  
    & & & \multicolumn{2}{p{0.3\linewidth}}{
        \cellcolor[rgb]{0.9,0.9,0.9}{
            \makecell[{{p{\linewidth}}}]{
                \texttt{\tiny{[GM$|$GM]}}
                {'hidden': {'team': ['diena', 'šūvis', 'stulpas', 'uoga', 'filmas', 'avarija', 'molis', 'obuolys', 'automobilis'], 'innocent': ['šakutė', 'skalbimo mašina', 'burna', 'citrina', 'amazonė', 'sūpynės', 'vamzdis', 'kengūra'], 'opponent': ['kojinė', 'šaltas', 'baterija', 'romanas', 'krūtinė', 'afrika', 'naktis', 'meksika'], 'assassin': []}, 'revealed': {'team': {'team': [], 'innocent': [], 'opponent': [], 'assassin': []}, 'opponent': {'team': [], 'innocent': [], 'opponent': [], 'assassin': []}}}
            }
        }
    }
    & & \\ \\

    \theutterance \stepcounter{utterance}  
    & & \multicolumn{4}{p{0.6\linewidth}}{
        \cellcolor[rgb]{0.9,0.9,0.9}{
            \makecell[{{p{\linewidth}}}]{
                \texttt{\tiny{[P1$\langle$GM]}}
                \texttt{Pateikite vieną vieno žodžio užuominą, susijusią su kai kuriais jūsų komandos žodžiais iš toliau pateikto sąrašo. Galite pasirinkti, su kiek ir kokiais žodžiais bus susijusi jūsų užuomina.} \\
\texttt{Tikslas {-} kad jūsų partneris, remdamasis jūsų pateiktomis užuominomis, atspėtų visus jūsų komandos žodžius, kol kita komanda neatskleidė visų savo žodžių. Jei jūsų partneris atspėja "priešininko" arba "atitraukiančio dėmesį" žodį, jo spėjimo eilė baigiasi ir priešininkų komanda gali atlikti savo ėjimą. Jei jūsų partneris atspėja "žudantį" žodį, žaidimas baigiasi ir jūsų komanda iš karto pralaimi.} \\
\texttt{Kuo daugiau žodžių, su kuriais susijusi jūsų užuomina, tuo greičiau galite laimėti žaidimą, tačiau partneriui gali būti sunkiau atspėti asociacijas, todėl išmintingai pasirinkite užuominą. Užuominos žodis turi būti semantiškai susijęs su tiksliniais žodžiais, jis negali būti vienas iš sąrašuose esančių žodžių arba sudaryti jų dalių.} \\
\\ 
\texttt{Visada pateikite vieno žodžio užuominą ir kableliais atskirtą susijusių tikslinių žodžių sąrašą toliau nurodyta forma ir atsakymus pateikite kuo trumpesnius, niekada nepateikite jokio kito teksto, nei reikalaujama šioje formoje:} \\
\\ 
\texttt{UŽUOMINA: <ŽODIS>} \\
\texttt{TIKSLAI: <TIKSLAI>} \\
\\ 
\texttt{Jūsų komandos žodžiai yra: diena, šūvis, stulpas, uoga, filmas, avarija, molis, obuolys, automobilis.} \\
\texttt{Jūsų priešininko komandos žodžiai yra: kojinė, šaltas, baterija, romanas, krūtinė, afrika, naktis, meksika.} \\
\texttt{Atitraukimo žodžiai yra: šakutė, skalbimo mašina, burna, citrina, amazonė, sūpynės, vamzdis, kengūra.} \\
\texttt{Žudymo žodžiai yra: .} \\
            }
        }
    }
    & \\ \\

    \theutterance \stepcounter{utterance}  
    & \multicolumn{4}{p{0.6\linewidth}}{
        \cellcolor[rgb]{0.9,0.9,0.9}{
            \makecell[{{p{\linewidth}}}]{
                \texttt{\tiny{[P1$\rangle$GM]}}
                \texttt{UŽUOMINA: vaisius} \\
\texttt{TIKSLAI: uoga, obuolys} \\
            }
        }
    }
    & & \\ \\

    \theutterance \stepcounter{utterance}  
    & & & \multicolumn{2}{p{0.3\linewidth}}{
        \cellcolor[rgb]{0.9,0.9,0.9}{
            \makecell[{{p{\linewidth}}}]{
                \texttt{\tiny{[GM$|$GM]}}
                \texttt{vaisius} \\
            }
        }
    }
    & & \\ \\

    \theutterance \stepcounter{utterance}  
    & & & \multicolumn{2}{p{0.3\linewidth}}{
        \cellcolor[rgb]{0.9,0.9,0.9}{
            \makecell[{{p{\linewidth}}}]{
                \texttt{\tiny{[GM$|$GM]}}
                ['uoga', 'obuolys']
            }
        }
    }
    & & \\ \\

    \theutterance \stepcounter{utterance}  
    & & & \multicolumn{2}{p{0.3\linewidth}}{
        \cellcolor[rgb]{0.9,0.9,0.9}{
            \makecell[{{p{\linewidth}}}]{
                \texttt{\tiny{[GM$|$GM]}}
                {'word': 'uoga', 'assignment': 'team'}
            }
        }
    }
    & & \\ \\

    \theutterance \stepcounter{utterance}  
    & & & \multicolumn{2}{p{0.3\linewidth}}{
        \cellcolor[rgb]{0.9,0.9,0.9}{
            \makecell[{{p{\linewidth}}}]{
                \texttt{\tiny{[GM$|$GM]}}
                {'word': 'obuolys', 'assignment': 'team'}
            }
        }
    }
    & & \\ \\

    \theutterance \stepcounter{utterance}  
    & & \multicolumn{4}{p{0.6\linewidth}}{
        \cellcolor[rgb]{0.9,0.9,0.9}{
            \makecell[{{p{\linewidth}}}]{
                \texttt{\tiny{[GM$\rangle$P2]}}
                \texttt{Pateikite kableliais atskirtą sąrašą iki 2 žodžių iš toliau pateikto sąrašo, kurie geriausiai siejasi arba yra labiausiai susiję su žodžiu 'vaisius'. Visada pradėkite savo spėjimų sąrašą žodžiais SPĖJIMAS:  ir į atsakymą neįtraukite jokio kito teksto.} \\
\\ 
\texttt{avarija, šaltas, stulpas, burna, citrina, skalbimo mašina, romanas, baterija, amazonė, krūtinė, naktis, diena, obuolys, filmas, kengūra, sūpynės, šakutė, molis, kojinė, meksika, afrika, automobilis, šūvis, uoga, vamzdis} \\
            }
        }
    }
    & \\ \\

    \theutterance \stepcounter{utterance}  
    & & & \multicolumn{4}{p{0.6\linewidth}}{
        \cellcolor[rgb]{0.9,0.9,0.9}{
            \makecell[{{p{\linewidth}}}]{
                \texttt{\tiny{[GM$\langle$P2]}}
                \texttt{SPĖJIMAS: obuolys, citrina} \\
            }
        }
    }
     \\ \\

    \theutterance \stepcounter{utterance}  
    & & & \multicolumn{2}{p{0.3\linewidth}}{
        \cellcolor[rgb]{0.9,0.9,0.9}{
            \makecell[{{p{\linewidth}}}]{
                \texttt{\tiny{[GM$|$GM]}}
                ['obuolys', 'citrina']
            }
        }
    }
    & & \\ \\

    \theutterance \stepcounter{utterance}  
    & & & \multicolumn{2}{p{0.3\linewidth}}{
        \cellcolor[rgb]{0.9,0.9,0.9}{
            \makecell[{{p{\linewidth}}}]{
                \texttt{\tiny{[GM$|$GM]}}
                {'word': 'obuolys', 'assignment': 'team'}
            }
        }
    }
    & & \\ \\

    \theutterance \stepcounter{utterance}  
    & & & \multicolumn{2}{p{0.3\linewidth}}{
        \cellcolor[rgb]{0.9,0.9,0.9}{
            \makecell[{{p{\linewidth}}}]{
                \texttt{\tiny{[GM$|$GM]}}
                {'word': 'obuolys', 'assignment': 'team'}
            }
        }
    }
    & & \\ \\

    \theutterance \stepcounter{utterance}  
    & & & \multicolumn{2}{p{0.3\linewidth}}{
        \cellcolor[rgb]{0.9,0.9,0.9}{
            \makecell[{{p{\linewidth}}}]{
                \texttt{\tiny{[GM$|$GM]}}
                {'word': 'citrina', 'assignment': 'innocent'}
            }
        }
    }
    & & \\ \\

    \theutterance \stepcounter{utterance}  
    & & & \multicolumn{2}{p{0.3\linewidth}}{
        \cellcolor[rgb]{0.9,0.9,0.9}{
            \makecell[{{p{\linewidth}}}]{
                \texttt{\tiny{[GM$|$GM]}}
                \texttt{citrina} \\
            }
        }
    }
    & & \\ \\

    \theutterance \stepcounter{utterance}  
    & & & \multicolumn{2}{p{0.3\linewidth}}{
        \cellcolor[rgb]{0.9,0.9,0.9}{
            \makecell[{{p{\linewidth}}}]{
                \texttt{\tiny{[GM$|$GM]}}
                {'word': 'romanas', 'assignment': 'opponent'}
            }
        }
    }
    & & \\ \\

    \theutterance \stepcounter{utterance}  
    & & & \multicolumn{2}{p{0.3\linewidth}}{
        \cellcolor[rgb]{0.9,0.9,0.9}{
            \makecell[{{p{\linewidth}}}]{
                \texttt{\tiny{[GM$|$GM]}}
                {'hidden': {'team': ['diena', 'šūvis', 'stulpas', 'uoga', 'filmas', 'avarija', 'molis', 'automobilis'], 'innocent': ['šakutė', 'skalbimo mašina', 'burna', 'amazonė', 'sūpynės', 'vamzdis', 'kengūra'], 'opponent': ['kojinė', 'šaltas', 'baterija', 'krūtinė', 'afrika', 'naktis', 'meksika'], 'assassin': []}, 'revealed': {'team': {'team': ['obuolys'], 'innocent': ['citrina'], 'opponent': [], 'assassin': []}, 'opponent': {'team': [], 'innocent': [], 'opponent': ['romanas'], 'assassin': []}}}
            }
        }
    }
    & & \\ \\

    \theutterance \stepcounter{utterance}  
    & & \multicolumn{4}{p{0.6\linewidth}}{
        \cellcolor[rgb]{0.9,0.9,0.9}{
            \makecell[{{p{\linewidth}}}]{
                \texttt{\tiny{[P1$\langle$GM]}}
                \texttt{Žodis obuolys buvo atspėtas teisingai. Žodis citrina buvo atspėtas, tačiau tai nekaltas žodis. Jūsų komandos draugo ėjimas baigėsi.} \\
\texttt{Dabar pateikite kitą užuominą, susijusią su kai kuriais likusiais jūsų komandos žodžiais, ir susijusių tikslinių žodžių sąrašą. Nepamirškite užuominą pradėti žodžiais "UŽUOMINA: ", įrašykite naują eilutę, o tikslinių žodžių sąrašą, atskirtą kableliais, pradėkite žodžiais "TIKSLAI: ". Pastaba: palyginti su ankstesniais prašymais, kai kurie žodžiai iš sąrašų buvo pašalinti.} \\
\\ 
\texttt{Likę jūsų komandos žodžiai yra šie: diena, šūvis, stulpas, uoga, filmas, avarija, molis, automobilis.} \\
\texttt{Likę priešininko žodžiai yra šie: kojinė, šaltas, baterija, romanas, krūtinė, afrika, naktis, meksika.} \\
\texttt{Likę atitraukiantys žodžiai yra: šakutė, skalbimo mašina, burna, amazonė, sūpynės, vamzdis, kengūra.} \\
\texttt{Likę nužudymo žodžiai yra: .} \\
            }
        }
    }
    & \\ \\

    \theutterance \stepcounter{utterance}  
    & \multicolumn{4}{p{0.6\linewidth}}{
        \cellcolor[rgb]{0.9,0.9,0.9}{
            \makecell[{{p{\linewidth}}}]{
                \texttt{\tiny{[P1$\rangle$GM]}}
                \texttt{UŽUOMINA: susidūrimas} \\
\texttt{TIKSLAI: avarija, automobilis} \\
            }
        }
    }
    & & \\ \\

    \theutterance \stepcounter{utterance}  
    & & & \multicolumn{2}{p{0.3\linewidth}}{
        \cellcolor[rgb]{0.9,0.9,0.9}{
            \makecell[{{p{\linewidth}}}]{
                \texttt{\tiny{[GM$|$GM]}}
                \texttt{susidūrimas} \\
            }
        }
    }
    & & \\ \\

    \theutterance \stepcounter{utterance}  
    & & & \multicolumn{2}{p{0.3\linewidth}}{
        \cellcolor[rgb]{0.9,0.9,0.9}{
            \makecell[{{p{\linewidth}}}]{
                \texttt{\tiny{[GM$|$GM]}}
                ['avarija', 'automobilis']
            }
        }
    }
    & & \\ \\

    \theutterance \stepcounter{utterance}  
    & & & \multicolumn{2}{p{0.3\linewidth}}{
        \cellcolor[rgb]{0.9,0.9,0.9}{
            \makecell[{{p{\linewidth}}}]{
                \texttt{\tiny{[GM$|$GM]}}
                {'word': 'avarija', 'assignment': 'team'}
            }
        }
    }
    & & \\ \\

    \theutterance \stepcounter{utterance}  
    & & & \multicolumn{2}{p{0.3\linewidth}}{
        \cellcolor[rgb]{0.9,0.9,0.9}{
            \makecell[{{p{\linewidth}}}]{
                \texttt{\tiny{[GM$|$GM]}}
                {'word': 'automobilis', 'assignment': 'team'}
            }
        }
    }
    & & \\ \\

    \theutterance \stepcounter{utterance}  
    & & \multicolumn{4}{p{0.6\linewidth}}{
        \cellcolor[rgb]{0.9,0.9,0.9}{
            \makecell[{{p{\linewidth}}}]{
                \texttt{\tiny{[GM$\rangle$P2]}}
                \texttt{Žodis obuolys buvo atspėtas teisingai. Žodis citrina buvo atspėtas, tačiau tai nekaltas žodis. Jūsų ėjimas baigėsi.} \\
\texttt{Dabar pateikite dar vieną kableliais atskirtą sąrašą iš bent 1 ir ne daugiau kaip 2 žodžių iš toliau pateikto žodžių sąrašo, kurie geriausiai siejasi arba yra labiausiai susiję su žodžiu 'susidūrimas'. Nepamirškite savo atsakymą pradėti SPĖJIMAS: . Pastaba: palyginti su ankstesniais prašymais, kai kurie žodžiai iš sąrašo buvo pašalinti.} \\
\\ 
\texttt{avarija, šaltas, stulpas, burna, skalbimo mašina, baterija, amazonė, krūtinė, naktis, diena, filmas, kengūra, sūpynės, šakutė, molis, kojinė, meksika, afrika, automobilis, šūvis, uoga, vamzdis} \\
            }
        }
    }
    & \\ \\

    \theutterance \stepcounter{utterance}  
    & & & \multicolumn{4}{p{0.6\linewidth}}{
        \cellcolor[rgb]{0.9,0.9,0.9}{
            \makecell[{{p{\linewidth}}}]{
                \texttt{\tiny{[GM$\langle$P2]}}
                \texttt{SPĖJIMAS: avarija, automobilis} \\
            }
        }
    }
     \\ \\

    \theutterance \stepcounter{utterance}  
    & & & \multicolumn{2}{p{0.3\linewidth}}{
        \cellcolor[rgb]{0.9,0.9,0.9}{
            \makecell[{{p{\linewidth}}}]{
                \texttt{\tiny{[GM$|$GM]}}
                ['avarija', 'automobilis']
            }
        }
    }
    & & \\ \\

    \theutterance \stepcounter{utterance}  
    & & & \multicolumn{2}{p{0.3\linewidth}}{
        \cellcolor[rgb]{0.9,0.9,0.9}{
            \makecell[{{p{\linewidth}}}]{
                \texttt{\tiny{[GM$|$GM]}}
                {'word': 'avarija', 'assignment': 'team'}
            }
        }
    }
    & & \\ \\

    \theutterance \stepcounter{utterance}  
    & & & \multicolumn{2}{p{0.3\linewidth}}{
        \cellcolor[rgb]{0.9,0.9,0.9}{
            \makecell[{{p{\linewidth}}}]{
                \texttt{\tiny{[GM$|$GM]}}
                {'word': 'avarija', 'assignment': 'team'}
            }
        }
    }
    & & \\ \\

    \theutterance \stepcounter{utterance}  
    & & & \multicolumn{2}{p{0.3\linewidth}}{
        \cellcolor[rgb]{0.9,0.9,0.9}{
            \makecell[{{p{\linewidth}}}]{
                \texttt{\tiny{[GM$|$GM]}}
                {'word': 'automobilis', 'assignment': 'team'}
            }
        }
    }
    & & \\ \\

    \theutterance \stepcounter{utterance}  
    & & & \multicolumn{2}{p{0.3\linewidth}}{
        \cellcolor[rgb]{0.9,0.9,0.9}{
            \makecell[{{p{\linewidth}}}]{
                \texttt{\tiny{[GM$|$GM]}}
                {'word': 'automobilis', 'assignment': 'team'}
            }
        }
    }
    & & \\ \\

    \theutterance \stepcounter{utterance}  
    & & & \multicolumn{2}{p{0.3\linewidth}}{
        \cellcolor[rgb]{0.9,0.9,0.9}{
            \makecell[{{p{\linewidth}}}]{
                \texttt{\tiny{[GM$|$GM]}}
                {'word': 'šaltas', 'assignment': 'opponent'}
            }
        }
    }
    & & \\ \\

    \theutterance \stepcounter{utterance}  
    & & & \multicolumn{2}{p{0.3\linewidth}}{
        \cellcolor[rgb]{0.9,0.9,0.9}{
            \makecell[{{p{\linewidth}}}]{
                \texttt{\tiny{[GM$|$GM]}}
                {'hidden': {'team': ['diena', 'šūvis', 'stulpas', 'uoga', 'filmas', 'molis'], 'innocent': ['šakutė', 'skalbimo mašina', 'burna', 'amazonė', 'sūpynės', 'vamzdis', 'kengūra'], 'opponent': ['kojinė', 'baterija', 'krūtinė', 'afrika', 'naktis', 'meksika'], 'assassin': []}, 'revealed': {'team': {'team': ['obuolys', 'avarija', 'automobilis'], 'innocent': ['citrina'], 'opponent': [], 'assassin': []}, 'opponent': {'team': [], 'innocent': [], 'opponent': ['romanas', 'šaltas'], 'assassin': []}}}
            }
        }
    }
    & & \\ \\

    \theutterance \stepcounter{utterance}  
    & & \multicolumn{4}{p{0.6\linewidth}}{
        \cellcolor[rgb]{0.9,0.9,0.9}{
            \makecell[{{p{\linewidth}}}]{
                \texttt{\tiny{[P1$\langle$GM]}}
                \texttt{Žodžiai avarija, automobilis buvo atspėti teisingai. Jūsų komandos draugo ėjimas baigėsi.} \\
\texttt{Dabar pateikite kitą užuominą, susijusią su kai kuriais likusiais jūsų komandos žodžiais, ir susijusių tikslinių žodžių sąrašą. Nepamirškite užuominą pradėti žodžiais "UŽUOMINA: ", įrašykite naują eilutę, o tikslinių žodžių sąrašą, atskirtą kableliais, pradėkite žodžiais "TIKSLAI: ". Pastaba: palyginti su ankstesniais prašymais, kai kurie žodžiai iš sąrašų buvo pašalinti.} \\
\\ 
\texttt{Likę jūsų komandos žodžiai yra šie: diena, šūvis, stulpas, uoga, filmas, molis.} \\
\texttt{Likę priešininko žodžiai yra šie: kojinė, šaltas, baterija, krūtinė, afrika, naktis, meksika.} \\
\texttt{Likę atitraukiantys žodžiai yra: šakutė, skalbimo mašina, burna, amazonė, sūpynės, vamzdis, kengūra.} \\
\texttt{Likę nužudymo žodžiai yra: .} \\
            }
        }
    }
    & \\ \\

    \theutterance \stepcounter{utterance}  
    & \multicolumn{4}{p{0.6\linewidth}}{
        \cellcolor[rgb]{0.9,0.9,0.9}{
            \makecell[{{p{\linewidth}}}]{
                \texttt{\tiny{[P1$\rangle$GM]}}
                \texttt{UŽUOMINA: kinas} \\
\texttt{TIKSLAI: filmas} \\
            }
        }
    }
    & & \\ \\

    \theutterance \stepcounter{utterance}  
    & & & \multicolumn{2}{p{0.3\linewidth}}{
        \cellcolor[rgb]{0.9,0.9,0.9}{
            \makecell[{{p{\linewidth}}}]{
                \texttt{\tiny{[GM$|$GM]}}
                \texttt{kinas} \\
            }
        }
    }
    & & \\ \\

    \theutterance \stepcounter{utterance}  
    & & & \multicolumn{2}{p{0.3\linewidth}}{
        \cellcolor[rgb]{0.9,0.9,0.9}{
            \makecell[{{p{\linewidth}}}]{
                \texttt{\tiny{[GM$|$GM]}}
                ['filmas']
            }
        }
    }
    & & \\ \\

    \theutterance \stepcounter{utterance}  
    & & & \multicolumn{2}{p{0.3\linewidth}}{
        \cellcolor[rgb]{0.9,0.9,0.9}{
            \makecell[{{p{\linewidth}}}]{
                \texttt{\tiny{[GM$|$GM]}}
                {'word': 'filmas', 'assignment': 'team'}
            }
        }
    }
    & & \\ \\

    \theutterance \stepcounter{utterance}  
    & & \multicolumn{4}{p{0.6\linewidth}}{
        \cellcolor[rgb]{0.9,0.9,0.9}{
            \makecell[{{p{\linewidth}}}]{
                \texttt{\tiny{[GM$\rangle$P2]}}
                \texttt{Žodžiai avarija, automobilis buvo atspėti teisingai. Jūsų ėjimas baigėsi.} \\
\texttt{Dabar pateikite dar vieną kableliais atskirtą sąrašą iš bent 1 ir ne daugiau kaip 1 žodžių iš toliau pateikto žodžių sąrašo, kurie geriausiai siejasi arba yra labiausiai susiję su žodžiu 'kinas'. Nepamirškite savo atsakymą pradėti SPĖJIMAS: . Pastaba: palyginti su ankstesniais prašymais, kai kurie žodžiai iš sąrašo buvo pašalinti.} \\
\\ 
\texttt{stulpas, burna, skalbimo mašina, baterija, amazonė, krūtinė, naktis, diena, filmas, kengūra, sūpynės, šakutė, molis, kojinė, meksika, afrika, šūvis, uoga, vamzdis} \\
            }
        }
    }
    & \\ \\

    \theutterance \stepcounter{utterance}  
    & & & \multicolumn{4}{p{0.6\linewidth}}{
        \cellcolor[rgb]{0.9,0.9,0.9}{
            \makecell[{{p{\linewidth}}}]{
                \texttt{\tiny{[GM$\langle$P2]}}
                \texttt{SPĖJIMAS: filmas} \\
            }
        }
    }
     \\ \\

    \theutterance \stepcounter{utterance}  
    & & & \multicolumn{2}{p{0.3\linewidth}}{
        \cellcolor[rgb]{0.9,0.9,0.9}{
            \makecell[{{p{\linewidth}}}]{
                \texttt{\tiny{[GM$|$GM]}}
                ['filmas']
            }
        }
    }
    & & \\ \\

    \theutterance \stepcounter{utterance}  
    & & & \multicolumn{2}{p{0.3\linewidth}}{
        \cellcolor[rgb]{0.9,0.9,0.9}{
            \makecell[{{p{\linewidth}}}]{
                \texttt{\tiny{[GM$|$GM]}}
                {'word': 'filmas', 'assignment': 'team'}
            }
        }
    }
    & & \\ \\

    \theutterance \stepcounter{utterance}  
    & & & \multicolumn{2}{p{0.3\linewidth}}{
        \cellcolor[rgb]{0.9,0.9,0.9}{
            \makecell[{{p{\linewidth}}}]{
                \texttt{\tiny{[GM$|$GM]}}
                {'word': 'filmas', 'assignment': 'team'}
            }
        }
    }
    & & \\ \\

    \theutterance \stepcounter{utterance}  
    & & & \multicolumn{2}{p{0.3\linewidth}}{
        \cellcolor[rgb]{0.9,0.9,0.9}{
            \makecell[{{p{\linewidth}}}]{
                \texttt{\tiny{[GM$|$GM]}}
                {'word': 'baterija', 'assignment': 'opponent'}
            }
        }
    }
    & & \\ \\

    \theutterance \stepcounter{utterance}  
    & & & \multicolumn{2}{p{0.3\linewidth}}{
        \cellcolor[rgb]{0.9,0.9,0.9}{
            \makecell[{{p{\linewidth}}}]{
                \texttt{\tiny{[GM$|$GM]}}
                {'hidden': {'team': ['diena', 'šūvis', 'stulpas', 'uoga', 'molis'], 'innocent': ['šakutė', 'skalbimo mašina', 'burna', 'amazonė', 'sūpynės', 'vamzdis', 'kengūra'], 'opponent': ['kojinė', 'krūtinė', 'afrika', 'naktis', 'meksika'], 'assassin': []}, 'revealed': {'team': {'team': ['obuolys', 'avarija', 'automobilis', 'filmas'], 'innocent': ['citrina'], 'opponent': [], 'assassin': []}, 'opponent': {'team': [], 'innocent': [], 'opponent': ['romanas', 'šaltas', 'baterija'], 'assassin': []}}}
            }
        }
    }
    & & \\ \\

    \theutterance \stepcounter{utterance}  
    & & \multicolumn{4}{p{0.6\linewidth}}{
        \cellcolor[rgb]{0.9,0.9,0.9}{
            \makecell[{{p{\linewidth}}}]{
                \texttt{\tiny{[P1$\langle$GM]}}
                \texttt{Žodis filmas buvo atspėtas teisingai. Jūsų komandos draugo ėjimas baigėsi.} \\
\texttt{Dabar pateikite kitą užuominą, susijusią su kai kuriais likusiais jūsų komandos žodžiais, ir susijusių tikslinių žodžių sąrašą. Nepamirškite užuominą pradėti žodžiais "UŽUOMINA: ", įrašykite naują eilutę, o tikslinių žodžių sąrašą, atskirtą kableliais, pradėkite žodžiais "TIKSLAI: ". Pastaba: palyginti su ankstesniais prašymais, kai kurie žodžiai iš sąrašų buvo pašalinti.} \\
\\ 
\texttt{Likę jūsų komandos žodžiai yra šie: diena, šūvis, stulpas, uoga, molis.} \\
\texttt{Likę priešininko žodžiai yra šie: kojinė, baterija, krūtinė, afrika, naktis, meksika.} \\
\texttt{Likę atitraukiantys žodžiai yra: šakutė, skalbimo mašina, burna, amazonė, sūpynės, vamzdis, kengūra.} \\
\texttt{Likę nužudymo žodžiai yra: .} \\
            }
        }
    }
    & \\ \\

    \theutterance \stepcounter{utterance}  
    & \multicolumn{4}{p{0.6\linewidth}}{
        \cellcolor[rgb]{0.9,0.9,0.9}{
            \makecell[{{p{\linewidth}}}]{
                \texttt{\tiny{[P1$\rangle$GM]}}
                \texttt{UŽUOMINA: keramika} \\
\texttt{TIKSLAI: molis} \\
            }
        }
    }
    & & \\ \\

    \theutterance \stepcounter{utterance}  
    & & & \multicolumn{2}{p{0.3\linewidth}}{
        \cellcolor[rgb]{0.9,0.9,0.9}{
            \makecell[{{p{\linewidth}}}]{
                \texttt{\tiny{[GM$|$GM]}}
                \texttt{keramika} \\
            }
        }
    }
    & & \\ \\

    \theutterance \stepcounter{utterance}  
    & & & \multicolumn{2}{p{0.3\linewidth}}{
        \cellcolor[rgb]{0.9,0.9,0.9}{
            \makecell[{{p{\linewidth}}}]{
                \texttt{\tiny{[GM$|$GM]}}
                ['molis']
            }
        }
    }
    & & \\ \\

    \theutterance \stepcounter{utterance}  
    & & & \multicolumn{2}{p{0.3\linewidth}}{
        \cellcolor[rgb]{0.9,0.9,0.9}{
            \makecell[{{p{\linewidth}}}]{
                \texttt{\tiny{[GM$|$GM]}}
                {'word': 'molis', 'assignment': 'team'}
            }
        }
    }
    & & \\ \\

    \theutterance \stepcounter{utterance}  
    & & \multicolumn{4}{p{0.6\linewidth}}{
        \cellcolor[rgb]{0.9,0.9,0.9}{
            \makecell[{{p{\linewidth}}}]{
                \texttt{\tiny{[GM$\rangle$P2]}}
                \texttt{Žodis filmas buvo atspėtas teisingai. Jūsų ėjimas baigėsi.} \\
\texttt{Dabar pateikite dar vieną kableliais atskirtą sąrašą iš bent 1 ir ne daugiau kaip 1 žodžių iš toliau pateikto žodžių sąrašo, kurie geriausiai siejasi arba yra labiausiai susiję su žodžiu 'keramika'. Nepamirškite savo atsakymą pradėti SPĖJIMAS: . Pastaba: palyginti su ankstesniais prašymais, kai kurie žodžiai iš sąrašo buvo pašalinti.} \\
\\ 
\texttt{stulpas, burna, skalbimo mašina, amazonė, krūtinė, naktis, diena, kengūra, sūpynės, šakutė, molis, kojinė, meksika, afrika, šūvis, uoga, vamzdis} \\
            }
        }
    }
    & \\ \\

    \theutterance \stepcounter{utterance}  
    & & & \multicolumn{4}{p{0.6\linewidth}}{
        \cellcolor[rgb]{0.9,0.9,0.9}{
            \makecell[{{p{\linewidth}}}]{
                \texttt{\tiny{[GM$\langle$P2]}}
                \texttt{SPĖJIMAS: molis} \\
            }
        }
    }
     \\ \\

    \theutterance \stepcounter{utterance}  
    & & & \multicolumn{2}{p{0.3\linewidth}}{
        \cellcolor[rgb]{0.9,0.9,0.9}{
            \makecell[{{p{\linewidth}}}]{
                \texttt{\tiny{[GM$|$GM]}}
                ['molis']
            }
        }
    }
    & & \\ \\

    \theutterance \stepcounter{utterance}  
    & & & \multicolumn{2}{p{0.3\linewidth}}{
        \cellcolor[rgb]{0.9,0.9,0.9}{
            \makecell[{{p{\linewidth}}}]{
                \texttt{\tiny{[GM$|$GM]}}
                {'word': 'molis', 'assignment': 'team'}
            }
        }
    }
    & & \\ \\

    \theutterance \stepcounter{utterance}  
    & & & \multicolumn{2}{p{0.3\linewidth}}{
        \cellcolor[rgb]{0.9,0.9,0.9}{
            \makecell[{{p{\linewidth}}}]{
                \texttt{\tiny{[GM$|$GM]}}
                {'word': 'molis', 'assignment': 'team'}
            }
        }
    }
    & & \\ \\

    \theutterance \stepcounter{utterance}  
    & & & \multicolumn{2}{p{0.3\linewidth}}{
        \cellcolor[rgb]{0.9,0.9,0.9}{
            \makecell[{{p{\linewidth}}}]{
                \texttt{\tiny{[GM$|$GM]}}
                {'word': 'kojinė', 'assignment': 'opponent'}
            }
        }
    }
    & & \\ \\

    \theutterance \stepcounter{utterance}  
    & & & \multicolumn{2}{p{0.3\linewidth}}{
        \cellcolor[rgb]{0.9,0.9,0.9}{
            \makecell[{{p{\linewidth}}}]{
                \texttt{\tiny{[GM$|$GM]}}
                {'hidden': {'team': ['diena', 'šūvis', 'stulpas', 'uoga'], 'innocent': ['šakutė', 'skalbimo mašina', 'burna', 'amazonė', 'sūpynės', 'vamzdis', 'kengūra'], 'opponent': ['krūtinė', 'afrika', 'naktis', 'meksika'], 'assassin': []}, 'revealed': {'team': {'team': ['obuolys', 'avarija', 'automobilis', 'filmas', 'molis'], 'innocent': ['citrina'], 'opponent': [], 'assassin': []}, 'opponent': {'team': [], 'innocent': [], 'opponent': ['romanas', 'šaltas', 'baterija', 'kojinė'], 'assassin': []}}}
            }
        }
    }
    & & \\ \\

    \theutterance \stepcounter{utterance}  
    & & \multicolumn{4}{p{0.6\linewidth}}{
        \cellcolor[rgb]{0.9,0.9,0.9}{
            \makecell[{{p{\linewidth}}}]{
                \texttt{\tiny{[P1$\langle$GM]}}
                \texttt{Žodis molis buvo atspėtas teisingai. Jūsų komandos draugo ėjimas baigėsi.} \\
\texttt{Dabar pateikite kitą užuominą, susijusią su kai kuriais likusiais jūsų komandos žodžiais, ir susijusių tikslinių žodžių sąrašą. Nepamirškite užuominą pradėti žodžiais "UŽUOMINA: ", įrašykite naują eilutę, o tikslinių žodžių sąrašą, atskirtą kableliais, pradėkite žodžiais "TIKSLAI: ". Pastaba: palyginti su ankstesniais prašymais, kai kurie žodžiai iš sąrašų buvo pašalinti.} \\
\\ 
\texttt{Likę jūsų komandos žodžiai yra šie: diena, šūvis, stulpas, uoga.} \\
\texttt{Likę priešininko žodžiai yra šie: kojinė, krūtinė, afrika, naktis, meksika.} \\
\texttt{Likę atitraukiantys žodžiai yra: šakutė, skalbimo mašina, burna, amazonė, sūpynės, vamzdis, kengūra.} \\
\texttt{Likę nužudymo žodžiai yra: .} \\
            }
        }
    }
    & \\ \\

    \theutterance \stepcounter{utterance}  
    & \multicolumn{4}{p{0.6\linewidth}}{
        \cellcolor[rgb]{0.9,0.9,0.9}{
            \makecell[{{p{\linewidth}}}]{
                \texttt{\tiny{[P1$\rangle$GM]}}
                \texttt{UŽUOMINA: strėlė} \\
\texttt{TIKSLAI: šūvis, stulpas, uoga} \\
            }
        }
    }
    & & \\ \\

    \theutterance \stepcounter{utterance}  
    & & & \multicolumn{2}{p{0.3\linewidth}}{
        \cellcolor[rgb]{0.9,0.9,0.9}{
            \makecell[{{p{\linewidth}}}]{
                \texttt{\tiny{[GM$|$GM]}}
                \texttt{strėlė} \\
            }
        }
    }
    & & \\ \\

    \theutterance \stepcounter{utterance}  
    & & & \multicolumn{2}{p{0.3\linewidth}}{
        \cellcolor[rgb]{0.9,0.9,0.9}{
            \makecell[{{p{\linewidth}}}]{
                \texttt{\tiny{[GM$|$GM]}}
                ['šūvis', 'stulpas', 'uoga']
            }
        }
    }
    & & \\ \\

    \theutterance \stepcounter{utterance}  
    & & & \multicolumn{2}{p{0.3\linewidth}}{
        \cellcolor[rgb]{0.9,0.9,0.9}{
            \makecell[{{p{\linewidth}}}]{
                \texttt{\tiny{[GM$|$GM]}}
                {'word': 'šūvis', 'assignment': 'team'}
            }
        }
    }
    & & \\ \\

    \theutterance \stepcounter{utterance}  
    & & & \multicolumn{2}{p{0.3\linewidth}}{
        \cellcolor[rgb]{0.9,0.9,0.9}{
            \makecell[{{p{\linewidth}}}]{
                \texttt{\tiny{[GM$|$GM]}}
                {'word': 'stulpas', 'assignment': 'team'}
            }
        }
    }
    & & \\ \\

    \theutterance \stepcounter{utterance}  
    & & & \multicolumn{2}{p{0.3\linewidth}}{
        \cellcolor[rgb]{0.9,0.9,0.9}{
            \makecell[{{p{\linewidth}}}]{
                \texttt{\tiny{[GM$|$GM]}}
                {'word': 'uoga', 'assignment': 'team'}
            }
        }
    }
    & & \\ \\

    \theutterance \stepcounter{utterance}  
    & & \multicolumn{4}{p{0.6\linewidth}}{
        \cellcolor[rgb]{0.9,0.9,0.9}{
            \makecell[{{p{\linewidth}}}]{
                \texttt{\tiny{[GM$\rangle$P2]}}
                \texttt{Žodis molis buvo atspėtas teisingai. Jūsų ėjimas baigėsi.} \\
\texttt{Dabar pateikite dar vieną kableliais atskirtą sąrašą iš bent 1 ir ne daugiau kaip 3 žodžių iš toliau pateikto žodžių sąrašo, kurie geriausiai siejasi arba yra labiausiai susiję su žodžiu 'strėlė'. Nepamirškite savo atsakymą pradėti SPĖJIMAS: . Pastaba: palyginti su ankstesniais prašymais, kai kurie žodžiai iš sąrašo buvo pašalinti.} \\
\\ 
\texttt{stulpas, burna, skalbimo mašina, amazonė, krūtinė, naktis, diena, kengūra, sūpynės, šakutė, meksika, afrika, šūvis, uoga, vamzdis} \\
            }
        }
    }
    & \\ \\

    \theutterance \stepcounter{utterance}  
    & & & \multicolumn{4}{p{0.6\linewidth}}{
        \cellcolor[rgb]{0.9,0.9,0.9}{
            \makecell[{{p{\linewidth}}}]{
                \texttt{\tiny{[GM$\langle$P2]}}
                \texttt{SPĖJIMAS: šūvis, amazonė, stulpas} \\
            }
        }
    }
     \\ \\

    \theutterance \stepcounter{utterance}  
    & & & \multicolumn{2}{p{0.3\linewidth}}{
        \cellcolor[rgb]{0.9,0.9,0.9}{
            \makecell[{{p{\linewidth}}}]{
                \texttt{\tiny{[GM$|$GM]}}
                ['šūvis', 'amazonė', 'stulpas']
            }
        }
    }
    & & \\ \\

    \theutterance \stepcounter{utterance}  
    & & & \multicolumn{2}{p{0.3\linewidth}}{
        \cellcolor[rgb]{0.9,0.9,0.9}{
            \makecell[{{p{\linewidth}}}]{
                \texttt{\tiny{[GM$|$GM]}}
                {'word': 'šūvis', 'assignment': 'team'}
            }
        }
    }
    & & \\ \\

    \theutterance \stepcounter{utterance}  
    & & & \multicolumn{2}{p{0.3\linewidth}}{
        \cellcolor[rgb]{0.9,0.9,0.9}{
            \makecell[{{p{\linewidth}}}]{
                \texttt{\tiny{[GM$|$GM]}}
                {'word': 'šūvis', 'assignment': 'team'}
            }
        }
    }
    & & \\ \\

    \theutterance \stepcounter{utterance}  
    & & & \multicolumn{2}{p{0.3\linewidth}}{
        \cellcolor[rgb]{0.9,0.9,0.9}{
            \makecell[{{p{\linewidth}}}]{
                \texttt{\tiny{[GM$|$GM]}}
                {'word': 'amazonė', 'assignment': 'innocent'}
            }
        }
    }
    & & \\ \\

    \theutterance \stepcounter{utterance}  
    & & & \multicolumn{2}{p{0.3\linewidth}}{
        \cellcolor[rgb]{0.9,0.9,0.9}{
            \makecell[{{p{\linewidth}}}]{
                \texttt{\tiny{[GM$|$GM]}}
                \texttt{amazonė} \\
            }
        }
    }
    & & \\ \\

    \theutterance \stepcounter{utterance}  
    & & & \multicolumn{2}{p{0.3\linewidth}}{
        \cellcolor[rgb]{0.9,0.9,0.9}{
            \makecell[{{p{\linewidth}}}]{
                \texttt{\tiny{[GM$|$GM]}}
                {'word': 'krūtinė', 'assignment': 'opponent'}
            }
        }
    }
    & & \\ \\

    \theutterance \stepcounter{utterance}  
    & & & \multicolumn{2}{p{0.3\linewidth}}{
        \cellcolor[rgb]{0.9,0.9,0.9}{
            \makecell[{{p{\linewidth}}}]{
                \texttt{\tiny{[GM$|$GM]}}
                {'hidden': {'team': ['diena', 'stulpas', 'uoga'], 'innocent': ['šakutė', 'skalbimo mašina', 'burna', 'sūpynės', 'vamzdis', 'kengūra'], 'opponent': ['afrika', 'naktis', 'meksika'], 'assassin': []}, 'revealed': {'team': {'team': ['obuolys', 'avarija', 'automobilis', 'filmas', 'molis', 'šūvis'], 'innocent': ['citrina', 'amazonė'], 'opponent': [], 'assassin': []}, 'opponent': {'team': [], 'innocent': [], 'opponent': ['romanas', 'šaltas', 'baterija', 'kojinė', 'krūtinė'], 'assassin': []}}}
            }
        }
    }
    & & \\ \\

    \theutterance \stepcounter{utterance}  
    & & \multicolumn{4}{p{0.6\linewidth}}{
        \cellcolor[rgb]{0.9,0.9,0.9}{
            \makecell[{{p{\linewidth}}}]{
                \texttt{\tiny{[P1$\langle$GM]}}
                \texttt{Žodis šūvis buvo atspėtas teisingai. Žodis amazonė buvo atspėtas, tačiau tai nekaltas žodis. Jūsų komandos draugo ėjimas baigėsi.} \\
\texttt{Dabar pateikite kitą užuominą, susijusią su kai kuriais likusiais jūsų komandos žodžiais, ir susijusių tikslinių žodžių sąrašą. Nepamirškite užuominą pradėti žodžiais "UŽUOMINA: ", įrašykite naują eilutę, o tikslinių žodžių sąrašą, atskirtą kableliais, pradėkite žodžiais "TIKSLAI: ". Pastaba: palyginti su ankstesniais prašymais, kai kurie žodžiai iš sąrašų buvo pašalinti.} \\
\\ 
\texttt{Likę jūsų komandos žodžiai yra šie: diena, stulpas, uoga.} \\
\texttt{Likę priešininko žodžiai yra šie: krūtinė, afrika, naktis, meksika.} \\
\texttt{Likę atitraukiantys žodžiai yra: šakutė, skalbimo mašina, burna, sūpynės, vamzdis, kengūra.} \\
\texttt{Likę nužudymo žodžiai yra: .} \\
            }
        }
    }
    & \\ \\

    \theutterance \stepcounter{utterance}  
    & \multicolumn{4}{p{0.6\linewidth}}{
        \cellcolor[rgb]{0.9,0.9,0.9}{
            \makecell[{{p{\linewidth}}}]{
                \texttt{\tiny{[P1$\rangle$GM]}}
                \texttt{UŽUOMINA: saulė} \\
\texttt{TIKSLAI: diena, uoga} \\
            }
        }
    }
    & & \\ \\

    \theutterance \stepcounter{utterance}  
    & & & \multicolumn{2}{p{0.3\linewidth}}{
        \cellcolor[rgb]{0.9,0.9,0.9}{
            \makecell[{{p{\linewidth}}}]{
                \texttt{\tiny{[GM$|$GM]}}
                \texttt{saulė} \\
            }
        }
    }
    & & \\ \\

    \theutterance \stepcounter{utterance}  
    & & & \multicolumn{2}{p{0.3\linewidth}}{
        \cellcolor[rgb]{0.9,0.9,0.9}{
            \makecell[{{p{\linewidth}}}]{
                \texttt{\tiny{[GM$|$GM]}}
                ['diena', 'uoga']
            }
        }
    }
    & & \\ \\

    \theutterance \stepcounter{utterance}  
    & & & \multicolumn{2}{p{0.3\linewidth}}{
        \cellcolor[rgb]{0.9,0.9,0.9}{
            \makecell[{{p{\linewidth}}}]{
                \texttt{\tiny{[GM$|$GM]}}
                {'word': 'diena', 'assignment': 'team'}
            }
        }
    }
    & & \\ \\

    \theutterance \stepcounter{utterance}  
    & & & \multicolumn{2}{p{0.3\linewidth}}{
        \cellcolor[rgb]{0.9,0.9,0.9}{
            \makecell[{{p{\linewidth}}}]{
                \texttt{\tiny{[GM$|$GM]}}
                {'word': 'uoga', 'assignment': 'team'}
            }
        }
    }
    & & \\ \\

    \theutterance \stepcounter{utterance}  
    & & \multicolumn{4}{p{0.6\linewidth}}{
        \cellcolor[rgb]{0.9,0.9,0.9}{
            \makecell[{{p{\linewidth}}}]{
                \texttt{\tiny{[GM$\rangle$P2]}}
                \texttt{Žodis šūvis buvo atspėtas teisingai. Žodis amazonė buvo atspėtas, tačiau tai nekaltas žodis. Jūsų ėjimas baigėsi.} \\
\texttt{Dabar pateikite dar vieną kableliais atskirtą sąrašą iš bent 1 ir ne daugiau kaip 2 žodžių iš toliau pateikto žodžių sąrašo, kurie geriausiai siejasi arba yra labiausiai susiję su žodžiu 'saulė'. Nepamirškite savo atsakymą pradėti SPĖJIMAS: . Pastaba: palyginti su ankstesniais prašymais, kai kurie žodžiai iš sąrašo buvo pašalinti.} \\
\\ 
\texttt{stulpas, burna, skalbimo mašina, naktis, diena, kengūra, sūpynės, šakutė, meksika, afrika, uoga, vamzdis} \\
            }
        }
    }
    & \\ \\

    \theutterance \stepcounter{utterance}  
    & & & \multicolumn{4}{p{0.6\linewidth}}{
        \cellcolor[rgb]{0.9,0.9,0.9}{
            \makecell[{{p{\linewidth}}}]{
                \texttt{\tiny{[GM$\langle$P2]}}
                \texttt{SPĖJIMAS: diena, naktis} \\
            }
        }
    }
     \\ \\

    \theutterance \stepcounter{utterance}  
    & & & \multicolumn{2}{p{0.3\linewidth}}{
        \cellcolor[rgb]{0.9,0.9,0.9}{
            \makecell[{{p{\linewidth}}}]{
                \texttt{\tiny{[GM$|$GM]}}
                ['diena', 'naktis']
            }
        }
    }
    & & \\ \\

    \theutterance \stepcounter{utterance}  
    & & & \multicolumn{2}{p{0.3\linewidth}}{
        \cellcolor[rgb]{0.9,0.9,0.9}{
            \makecell[{{p{\linewidth}}}]{
                \texttt{\tiny{[GM$|$GM]}}
                {'word': 'diena', 'assignment': 'team'}
            }
        }
    }
    & & \\ \\

    \theutterance \stepcounter{utterance}  
    & & & \multicolumn{2}{p{0.3\linewidth}}{
        \cellcolor[rgb]{0.9,0.9,0.9}{
            \makecell[{{p{\linewidth}}}]{
                \texttt{\tiny{[GM$|$GM]}}
                {'word': 'diena', 'assignment': 'team'}
            }
        }
    }
    & & \\ \\

    \theutterance \stepcounter{utterance}  
    & & & \multicolumn{2}{p{0.3\linewidth}}{
        \cellcolor[rgb]{0.9,0.9,0.9}{
            \makecell[{{p{\linewidth}}}]{
                \texttt{\tiny{[GM$|$GM]}}
                {'word': 'naktis', 'assignment': 'opponent'}
            }
        }
    }
    & & \\ \\

    \theutterance \stepcounter{utterance}  
    & & & \multicolumn{2}{p{0.3\linewidth}}{
        \cellcolor[rgb]{0.9,0.9,0.9}{
            \makecell[{{p{\linewidth}}}]{
                \texttt{\tiny{[GM$|$GM]}}
                \texttt{naktis} \\
            }
        }
    }
    & & \\ \\

    \theutterance \stepcounter{utterance}  
    & & & \multicolumn{2}{p{0.3\linewidth}}{
        \cellcolor[rgb]{0.9,0.9,0.9}{
            \makecell[{{p{\linewidth}}}]{
                \texttt{\tiny{[GM$|$GM]}}
                {'word': 'meksika', 'assignment': 'opponent'}
            }
        }
    }
    & & \\ \\

    \theutterance \stepcounter{utterance}  
    & & & \multicolumn{2}{p{0.3\linewidth}}{
        \cellcolor[rgb]{0.9,0.9,0.9}{
            \makecell[{{p{\linewidth}}}]{
                \texttt{\tiny{[GM$|$GM]}}
                {'hidden': {'team': ['stulpas', 'uoga'], 'innocent': ['šakutė', 'skalbimo mašina', 'burna', 'sūpynės', 'vamzdis', 'kengūra'], 'opponent': ['afrika'], 'assassin': []}, 'revealed': {'team': {'team': ['obuolys', 'avarija', 'automobilis', 'filmas', 'molis', 'šūvis', 'diena'], 'innocent': ['citrina', 'amazonė'], 'opponent': ['naktis'], 'assassin': []}, 'opponent': {'team': [], 'innocent': [], 'opponent': ['romanas', 'šaltas', 'baterija', 'kojinė', 'krūtinė', 'meksika'], 'assassin': []}}}
            }
        }
    }
    & & \\ \\

    \theutterance \stepcounter{utterance}  
    & & \multicolumn{4}{p{0.6\linewidth}}{
        \cellcolor[rgb]{0.9,0.9,0.9}{
            \makecell[{{p{\linewidth}}}]{
                \texttt{\tiny{[P1$\langle$GM]}}
                \texttt{Žodis diena buvo atspėtas teisingai. Žodis naktis buvo atspėtas, tačiau tai priešas žodis. Jūsų komandos draugo ėjimas baigėsi.} \\
\texttt{Dabar pateikite kitą užuominą, susijusią su kai kuriais likusiais jūsų komandos žodžiais, ir susijusių tikslinių žodžių sąrašą. Nepamirškite užuominą pradėti žodžiais "UŽUOMINA: ", įrašykite naują eilutę, o tikslinių žodžių sąrašą, atskirtą kableliais, pradėkite žodžiais "TIKSLAI: ". Pastaba: palyginti su ankstesniais prašymais, kai kurie žodžiai iš sąrašų buvo pašalinti.} \\
\\ 
\texttt{Likę jūsų komandos žodžiai yra šie: stulpas, uoga.} \\
\texttt{Likę priešininko žodžiai yra šie: afrika, meksika.} \\
\texttt{Likę atitraukiantys žodžiai yra: šakutė, skalbimo mašina, burna, sūpynės, vamzdis, kengūra.} \\
\texttt{Likę nužudymo žodžiai yra: .} \\
            }
        }
    }
    & \\ \\

    \theutterance \stepcounter{utterance}  
    & \multicolumn{4}{p{0.6\linewidth}}{
        \cellcolor[rgb]{0.9,0.9,0.9}{
            \makecell[{{p{\linewidth}}}]{
                \texttt{\tiny{[P1$\rangle$GM]}}
                \texttt{UŽUOMINA: kartis} \\
\texttt{TIKSLAI: stulpas} \\
            }
        }
    }
    & & \\ \\

    \theutterance \stepcounter{utterance}  
    & & & \multicolumn{2}{p{0.3\linewidth}}{
        \cellcolor[rgb]{0.9,0.9,0.9}{
            \makecell[{{p{\linewidth}}}]{
                \texttt{\tiny{[GM$|$GM]}}
                \texttt{kartis} \\
            }
        }
    }
    & & \\ \\

    \theutterance \stepcounter{utterance}  
    & & & \multicolumn{2}{p{0.3\linewidth}}{
        \cellcolor[rgb]{0.9,0.9,0.9}{
            \makecell[{{p{\linewidth}}}]{
                \texttt{\tiny{[GM$|$GM]}}
                ['stulpas']
            }
        }
    }
    & & \\ \\

    \theutterance \stepcounter{utterance}  
    & & & \multicolumn{2}{p{0.3\linewidth}}{
        \cellcolor[rgb]{0.9,0.9,0.9}{
            \makecell[{{p{\linewidth}}}]{
                \texttt{\tiny{[GM$|$GM]}}
                {'word': 'stulpas', 'assignment': 'team'}
            }
        }
    }
    & & \\ \\

    \theutterance \stepcounter{utterance}  
    & & \multicolumn{4}{p{0.6\linewidth}}{
        \cellcolor[rgb]{0.9,0.9,0.9}{
            \makecell[{{p{\linewidth}}}]{
                \texttt{\tiny{[GM$\rangle$P2]}}
                \texttt{Žodis diena buvo atspėtas teisingai. Žodis naktis buvo atspėtas, tačiau tai priešas žodis. Jūsų ėjimas baigėsi.} \\
\texttt{Dabar pateikite dar vieną kableliais atskirtą sąrašą iš bent 1 ir ne daugiau kaip 1 žodžių iš toliau pateikto žodžių sąrašo, kurie geriausiai siejasi arba yra labiausiai susiję su žodžiu 'kartis'. Nepamirškite savo atsakymą pradėti SPĖJIMAS: . Pastaba: palyginti su ankstesniais prašymais, kai kurie žodžiai iš sąrašo buvo pašalinti.} \\
\\ 
\texttt{stulpas, burna, skalbimo mašina, kengūra, sūpynės, šakutė, afrika, uoga, vamzdis} \\
            }
        }
    }
    & \\ \\

    \theutterance \stepcounter{utterance}  
    & & & \multicolumn{4}{p{0.6\linewidth}}{
        \cellcolor[rgb]{0.9,0.9,0.9}{
            \makecell[{{p{\linewidth}}}]{
                \texttt{\tiny{[GM$\langle$P2]}}
                \texttt{SPĖJIMAS: stulpas} \\
            }
        }
    }
     \\ \\

    \theutterance \stepcounter{utterance}  
    & & & \multicolumn{2}{p{0.3\linewidth}}{
        \cellcolor[rgb]{0.9,0.9,0.9}{
            \makecell[{{p{\linewidth}}}]{
                \texttt{\tiny{[GM$|$GM]}}
                ['stulpas']
            }
        }
    }
    & & \\ \\

    \theutterance \stepcounter{utterance}  
    & & & \multicolumn{2}{p{0.3\linewidth}}{
        \cellcolor[rgb]{0.9,0.9,0.9}{
            \makecell[{{p{\linewidth}}}]{
                \texttt{\tiny{[GM$|$GM]}}
                {'word': 'stulpas', 'assignment': 'team'}
            }
        }
    }
    & & \\ \\

    \theutterance \stepcounter{utterance}  
    & & & \multicolumn{2}{p{0.3\linewidth}}{
        \cellcolor[rgb]{0.9,0.9,0.9}{
            \makecell[{{p{\linewidth}}}]{
                \texttt{\tiny{[GM$|$GM]}}
                {'word': 'stulpas', 'assignment': 'team'}
            }
        }
    }
    & & \\ \\

    \theutterance \stepcounter{utterance}  
    & & & \multicolumn{2}{p{0.3\linewidth}}{
        \cellcolor[rgb]{0.9,0.9,0.9}{
            \makecell[{{p{\linewidth}}}]{
                \texttt{\tiny{[GM$|$GM]}}
                {'word': 'afrika', 'assignment': 'opponent'}
            }
        }
    }
    & & \\ \\

    \theutterance \stepcounter{utterance}  
    & & & \multicolumn{2}{p{0.3\linewidth}}{
        \cellcolor[rgb]{0.9,0.9,0.9}{
            \makecell[{{p{\linewidth}}}]{
                \texttt{\tiny{[GM$|$GM]}}
                {'hidden': {'team': ['uoga'], 'innocent': ['šakutė', 'skalbimo mašina', 'burna', 'sūpynės', 'vamzdis', 'kengūra'], 'opponent': [], 'assassin': []}, 'revealed': {'team': {'team': ['obuolys', 'avarija', 'automobilis', 'filmas', 'molis', 'šūvis', 'diena', 'stulpas'], 'innocent': ['citrina', 'amazonė'], 'opponent': ['naktis'], 'assassin': []}, 'opponent': {'team': [], 'innocent': [], 'opponent': ['romanas', 'šaltas', 'baterija', 'kojinė', 'krūtinė', 'meksika', 'afrika'], 'assassin': []}}}
            }
        }
    }
    & & \\ \\

    \theutterance \stepcounter{utterance}  
    & & \multicolumn{4}{p{0.6\linewidth}}{
        \cellcolor[rgb]{0.9,0.9,0.9}{
            \makecell[{{p{\linewidth}}}]{
                \texttt{\tiny{[P1$\langle$GM]}}
                \texttt{Žodis stulpas buvo atspėtas teisingai. Jūsų komandos draugo ėjimas baigėsi.} \\
\texttt{Dabar pateikite kitą užuominą, susijusią su kai kuriais likusiais jūsų komandos žodžiais, ir susijusių tikslinių žodžių sąrašą. Nepamirškite užuominą pradėti žodžiais "UŽUOMINA: ", įrašykite naują eilutę, o tikslinių žodžių sąrašą, atskirtą kableliais, pradėkite žodžiais "TIKSLAI: ". Pastaba: palyginti su ankstesniais prašymais, kai kurie žodžiai iš sąrašų buvo pašalinti.} \\
\\ 
\texttt{Likę jūsų komandos žodžiai yra šie: uoga.} \\
\texttt{Likę priešininko žodžiai yra šie: afrika.} \\
\texttt{Likę atitraukiantys žodžiai yra: šakutė, skalbimo mašina, burna, sūpynės, vamzdis, kengūra.} \\
\texttt{Likę nužudymo žodžiai yra: .} \\
            }
        }
    }
    & \\ \\

    \theutterance \stepcounter{utterance}  
    & \multicolumn{4}{p{0.6\linewidth}}{
        \cellcolor[rgb]{0.9,0.9,0.9}{
            \makecell[{{p{\linewidth}}}]{
                \texttt{\tiny{[P1$\rangle$GM]}}
                \texttt{UŽUOMINA: vaisius} \\
\texttt{TIKSLAI: uoga} \\
            }
        }
    }
    & & \\ \\

    \theutterance \stepcounter{utterance}  
    & & & \multicolumn{2}{p{0.3\linewidth}}{
        \cellcolor[rgb]{0.9,0.9,0.9}{
            \makecell[{{p{\linewidth}}}]{
                \texttt{\tiny{[GM$|$GM]}}
                \texttt{vaisius} \\
            }
        }
    }
    & & \\ \\

    \theutterance \stepcounter{utterance}  
    & & & \multicolumn{2}{p{0.3\linewidth}}{
        \cellcolor[rgb]{0.9,0.9,0.9}{
            \makecell[{{p{\linewidth}}}]{
                \texttt{\tiny{[GM$|$GM]}}
                ['uoga']
            }
        }
    }
    & & \\ \\

    \theutterance \stepcounter{utterance}  
    & & & \multicolumn{2}{p{0.3\linewidth}}{
        \cellcolor[rgb]{0.9,0.9,0.9}{
            \makecell[{{p{\linewidth}}}]{
                \texttt{\tiny{[GM$|$GM]}}
                {'word': 'uoga', 'assignment': 'team'}
            }
        }
    }
    & & \\ \\

    \theutterance \stepcounter{utterance}  
    & & & \multicolumn{2}{p{0.3\linewidth}}{
        \cellcolor[rgb]{0.9,0.9,0.9}{
            \makecell[{{p{\linewidth}}}]{
                \texttt{\tiny{[GM$|$GM]}}
                \texttt{opponent has won} \\
            }
        }
    }
    & & \\ \\

\end{supertabular}
}

\end{document}
